% OR03-S014 | lapisan asli: proyek kapstone masalah invers komposit
% segment-id: d90.orig.v1.tr03.seg0014
% capstone-id: d90.orig.v1.tr03.capstone.0001
% capstone-unit-id: d90.orig.v1.tr03.capstone.unit
% capstone-milestone-id: d90.orig.v1.tr03.capstone.milestone.0001
% capstone-milestone-id: d90.orig.v1.tr03.capstone.milestone.0002
% capstone-milestone-id: d90.orig.v1.tr03.capstone.milestone.0003
% capstone-milestone-id: d90.orig.v1.tr03.capstone.milestone.0004
% capstone-milestone-id: d90.orig.v1.tr03.capstone.milestone.0005
% capstone-milestone-id: d90.orig.v1.tr03.capstone.milestone.0006
% capstone-milestone-id: d90.orig.v1.tr03.capstone.milestone.0007
% capstone-hint-id: d90.orig.v1.tr03.capstone.hint.0001
% capstone-answer-id: d90.orig.v1.tr03.capstone.answer.0001
% capstone-solution-id: d90.orig.v1.tr03.capstone.solution.0001

\section{Proyek kapstone: masalah invers komposit yang tahan pencilan}
\label{orig03:capstone:invers-komposit}
\label{orig03:capstone:unit}

Proyek ini menyatukan pemodelan konveks, gradien Lipschitz, operator
proksimal, formulasi dual, pemisahan primal-dual, anggaran oracle, dan
sertifikat a posteriori. Instans tetap dibentuk oleh
\texttt{labs/original-03/kapstone-invers-komposit.py}: matriks
$A\in\mathbb{R}^{72\times48}$, sinyal renggang $x^\natural$, derau kecil,
dan delapan pencilan besar. Modelnya
\begin{equation}
 \min_x\; \Phi(x):=
 \sum_{i=1}^{72}h_\delta((Ax-b)_i)+\lambda\|x\|_1,
 \qquad
 h_\delta(r)=
 \begin{cases}
   \frac12r^2,&|r|\leq\delta,\\
   \delta(|r|-\frac12\delta),&|r|>\delta.
 \end{cases}
 \label{orig03:eq:capstone-model}
\end{equation}
FISTA dengan restart gradien dan PDHG dijalankan untuk jumlah iterasi yang sama.
Setiap iterasi inti masing-masing memakai satu aplikasi $A$ dan satu aplikasi
$A^\top$, sehingga perbandingan memakai anggaran operator yang sama. Jejak
CSV mencatat objektif, batas bawah dual layak, celah tersertifikasi, dan galat
rekonstruksi. SVG hanya merupakan visualisasi jejak tersebut.

\subsection*{Prompt}
\begin{exercise}[Tugas proyek dan tujuh tonggak]
Selesaikan tujuh tonggak berikut dan serahkan uraian matematis bersama JSON,
CSV, serta SVG yang dihasilkan tanpa perubahan konfigurasi.
\begin{enumerate}
  \item \label{orig03:capstone:milestone:0001}\textbf{Instans dan pencilan.} Rekonstruksi instans dari seed,
    verifikasi ukuran, support sinyal, dan delapan indeks pencilan.
  \item \label{orig03:capstone:milestone:0002}\textbf{Model konveks.} Buktikan kekonveksan \eqref{orig03:eq:capstone-model}
    dan jelaskan mengapa Huber lebih tahan pencilan daripada kuadrat murni.
  \item \label{orig03:capstone:milestone:0003}\textbf{Bagian mulus.} Turunkan gradien
    $A^\top\operatorname{clip}(Ax-b,-\delta,\delta)$ dan batas Lipschitz
    $L=\|A\|_2^2$.
  \item \label{orig03:capstone:milestone:0004}\textbf{Langkah proksimal.} Turunkan soft-threshold bagi
    $\lambda\|x\|_1$ dan tulis iterasi FISTA yang dipakai.
  \item \label{orig03:capstone:milestone:0005}\textbf{Dual dan PDHG.} Turunkan konjugat Huber, kendala dual, serta
    pembaruan PDHG; buktikan $\tau\sigma\|A\|_2^2<1$.
  \item \label{orig03:capstone:milestone:0006}\textbf{Anggaran sepadan.} Bandingkan kedua metode pada jumlah
    perkalian matriks-vektor yang sama, bukan pada waktu dinding.
  \item \label{orig03:capstone:milestone:0007}\textbf{Sertifikat dan interpretasi.} Verifikasi kelayakan dual,
    celah primal-dual, pemetaan gradien proksimal, dan perbaikan terhadap
    least squares; nyatakan juga apa yang tidak dibuktikan eksperimen ini.
\end{enumerate}
\end{exercise}

\subsection*{Petunjuk bertahap}
\label{orig03:capstone:hint:0001}
\begin{enumerate}
  \item Hash matriks dan data dalam JSON adalah identitas instans, bukan
    pengganti pemeriksaan ukuran dan seed.
  \item Turunan Huber adalah proyeksi skalar ke $[-\delta,\delta]$ dan
    bersifat satu-Lipschitz.
  \item Komposisi dengan $A$ mengalikan konstanta Lipschitz oleh
    $\|A\|_2^2$.
  \item Gunakan
    $\operatorname{prox}_{t\lambda\|\cdot\|_1}(v)
     =\operatorname{sign}(v)(|v|-t\lambda)_+$.
  \item Konjugat Huber adalah $\frac12u^2+\iota_{[-\delta,\delta]}(u)$;
    konjugat norma satu menghasilkan kendala
    $\|A^\top u\|_\infty\leq\lambda$.
  \item Kolom \texttt{matrix\_vector\_products} adalah sumbu perbandingan
    yang benar.
  \item Skala kandidat dual bila perlu untuk memperoleh batas bawah yang
    benar-benar layak sebelum menghitung celah.
\end{enumerate}

\subsection*{Jawaban singkat}
\label{orig03:capstone:answer:0001}
Model adalah jumlah fungsi konveks. Bagian Huber mulus memiliki gradien
$L$-Lipschitz dengan $L=\|A\|_2^2$, dan bagian norma satu memiliki proksimal
soft-threshold. Dualnya memaksimalkan
\[
 -b^\top u-\frac12\|u\|^2
 \quad\text{dengan}\quad
 \|u\|_\infty\leq\delta,
 \quad \|A^\top u\|_\infty\leq\lambda.
\]
FISTA dan PDHG mencapai objektif yang cocok pada anggaran operator sama;
batas dual layak memberi celah terukur, dan kedua rekonstruksi mengungguli
least squares pada instans berpencilan yang dibekukan.

\subsection*{Solusi acuan lengkap}
\label{orig03:capstone:solution:0001}
Karena $h_\delta$ dan norma satu konveks, komposisi afin dan penjumlahan
mempertahankan kekonveksan. Turunan Huber adalah
$h_\delta'(r)=\operatorname{clip}(r,-\delta,\delta)$. Proyeksi ini
satu-Lipschitz, sehingga
\[
 \|\nabla f(x)-\nabla f(y)\|
 \leq\|A\|_2^2\|x-y\|.
\]
Langkah $t<1/L$ yang dibekukan karena itu sah. Optimalisasi proksimal skalar
memberi soft-threshold, lalu ekstrapolasi Nesterov menghasilkan FISTA; restart
gradien memakai vektor yang sudah tersedia tanpa mengubah dua aplikasi operator inti per
iterasi.

Dari
$h_\delta^*(u)=\frac12u^2+\iota_{[-\delta,\delta]}(u)$ dan
$(\lambda\|\cdot\|_1)^*(v)=\iota_{\{\|v\|_\infty\leq\lambda\}}(v)$,
dual yang tertulis pada jawaban singkat mengikuti melalui Fenchel--Rockafellar.
Proksimal $h_\delta^*$ adalah pembagian oleh $1+\sigma$ diikuti kliping ke
$[-\delta,\delta]$; proksimal primal tetap soft-threshold. Pilihan
$\tau=\sigma=0.99/\|A\|_2$ memberi
$\tau\sigma\|A\|_2^2=0.99^2<1$.

Pada iterasi berhingga, kandidat dual dapat melanggar kendala kedua sedikit.
Penskalaan
\[
 \widehat u=\rho u,
 \qquad
 \rho=\min\left\{1,
 \frac{\lambda}{\|A^\top u\|_\infty}\right\}
\]
mempertahankan kendala kotak dan memastikan kelayakan. Nilai dual
$-b^\top\widehat u-\|\widehat u\|^2/2$ menjadi batas bawah rigor bagi
objektif primal. Program menguji kelayakan, celah terhadap batas terbaik,
norma pemetaan gradien proksimal FISTA, kesamaan anggaran, kedekatan kedua
objektif, dan galat rekonstruksi terhadap baseline least squares. Bukti ini
berlaku untuk instans dan toleransi yang dinyatakan; ia tidak mengklaim bahwa
parameter yang sama optimal bagi semua masalah invers atau semua pola
pencilan.

\paragraph{Hak dan firewall O018.}
Proyek, bukti, soal, solusi, kode, dan instans sintetik ini ditulis mandiri dan
tersedia berdasarkan CC BY-SA 4.0. Tidak ada materi O018 yang disalin,
diadaptasi, atau dijadikan templat tersembunyi; O018 tetap terpisah dari
lapisan penutupan Original-03.
